% !TEX root = ../main.tex
\clearpage
\phantomsection
\section*{3 \textbf{문헌 검토}}
\label{sec:literature-review}
\addcontentsline{toc}{section}{3 문헌 검토}

\phantomsection
\subsection*{3.1 그래프 기반 온체인 평판 점수 산출}
\label{sec:graph-based-on-chain-reputation-scoring}
\addcontentsline{toc}{subsection}{3.1 그래프 기반 온체인 평판 점수 산출}

선행 연구는 종종 이더리움 메인넷을 참조 체인으로 사용하지만, \textbf{본 학위논문의 주요 실증 설정은 Arbitrum One과 GMX V2 무기한 스왑 포지션 결과를 이용한 도메인 검증}이다. 아래 문헌은 그래프 기반 평판 방법의 배경으로서 여전히 관련이 있다.

온체인 평판에 대한 초기 연구는 그래프 이론, 특히 PageRank의 변형을 활용하여 지갑 주소를 ``신뢰성'' 순으로 순위화했다(Page et al., 1998). Do et al.(2023)은 이더리움 주소를 거래를 엣지로 하는 방향 그래프의 노드로 모델링한 뒤 PageRank 알고리즘을 적용하여 각 주소의 사회적 신용 점수를 계산한다. DeFi 성장은 무담보 대출 및 고객 보호를 위해 주체를 식별하고 순위화하는 방법을 요구한다는 것이 근거이다. 이 틀에서 주소는 다른 이들로부터 많거나 큰 거래로 ``신뢰''받을수록 더 높은 점수를 누적하며, 이는 중요한 웹 페이지가 많은 들어오는 링크를 갖는 것과 유사하다. PageRank는 따라서 평판의 대리변수로서 연결 강도와 활동을 측정한다. 후속 그래프 모델은 상호작용 빈도와 규모 등 여러 요인을 엣지 가중치에 결합한다(Do et al., 2023; Vathsala V. \& V. D. Sharma, 2025).


그래프 구성에서 핵심 고려는 엣지를 어떻게 가중할 것인가이다. 기본 구현은 각 거래를 동등하게 취급하지만, 후속 연구는 거래 가치를 가중치로 포함하여 큰 송금이 더 강한 신뢰를 나타낸다. HodgeRank 접근을 사용하여 Do et al.(2023)은 각 엣지에 이전된 ETH 양을 가중치로 두고 이러한 가중 흐름을 고려한 전역 순위를 구한다. 이는 거래량에 영향을 받는 이더리움 평판 순위를 산출하며, 고가치 상호작용이 당사자 간 더 큰 신뢰 또는 이해관계를 나타낸다는 직관에 기반한다. 또 다른 실무적 과제는 그래프 구조를 다루는 것이다: 이더리움 거래 그래프에는 PageRank 점수를 흡수하는 많은 고립 클러스터나 ``sink'' 노드가 포함될 수 있다. 전통적 PageRank와 같이 감쇠 인자를 사용하여 sink 및 작은 고립 부분그래프의 점수 팽창을 완화하고, 평판 점수가 고립된 활동이 아닌 더 넓은 네트워크 참여를 반영하도록 한다.


바닐라 PageRank, 즉 명시적 리스크 전파나 도메인별 조정 없이 거래 그래프에 적용된 수정되지 않은 PageRank 공식의 한계는 악의적이지만 활동적인 행위자와 진정으로 평판 있는 행위자를 구별하지 못할 수 있다는 것이다. Lin et al.(2024)은 거래 연결성만으로는 리스크 관련 구조를 놓칠 수 있으며, 리스크 신호가 네트워크를 통해 전파되어야 함을 보였다. 이는 정교화된 모델의 필요성을 강조한다. 제안된 개선에는 부정적 피드백 또는 리스크 레이블의 네트워크 전파 포함이 있다. 예를 들어 최근 연구는 이더리움 계정을 위한 네트워크 전파 기반 리스크 점수 모델 RiskProp을 도입했다(Lin et al., 2024). RiskProp은 주소에 대한 ``탈익명 점수''로 거래 그래프를 보강하고 방향 이분 그래프를 통해 리스크 신호를 전파한다. 이 접근은 거래 링크를 가로질러 리스크 지표를 전파하여 알려진 고위험 계정의 대다수를 플래그할 수 있었으며, 단순 연결성 측정이 놓칠 수 있는 의심 지갑을 효과적으로 식별했다. 이는 그래프 기반 방법이 순수 PageRank를 넘어 ``신용''이 금융 신뢰성 또는 규정 준수 리스크일 수 있는 온체인 신용 리스크 모델의 유사체로 확장될 수 있음을 보여준다.


그래프 기반 평판 시스템은 종종 시간적 역학을 고려한다는 점에 유의해야 한다. 실세계 신용도는 시간에 따라 변하므로, 일부 모델은 오래된 상호작용의 영향을 줄이기 위해 시간 감쇠 인자를 포함한다. Do et al.(2023)은 AWP에서 송금에 로지스틱 시간 감쇠를 적용하여 최근 활동이 오래된 엣지보다 더 강하게 기여하도록 한다. Lin et al.(2024)은 네트워크 상태가 진화할 때 전파 모델이 갱신되지 않으면 리스크 관련 구조가 가려질 수 있음을 보였다. 동적 평판 모델은 최근 거래에 더 높은 가중치를 부여하거나 주기적으로 오래된 엣지를 만료시킬 수 있다. 이는 주체가 과거의 긍정적 상호작용에만 의존하는 것을 방지하고, 나쁜 행동이 시간에 의해 희석되는 세탁(whitewashing) 문제를 다룬다. 그러나 적절한 균형이 중요하다: 과도한 감쇠는 오래 지속된 신뢰 관계를 무시할 수 있고, 너무 적은 감쇠는 오래된 정보가 점수를 왜곡할 수 있다.


\phantomsection
\subsection*{3.2 온체인 신용 리스크 예측 모델링}
\label{sec:predictive-modeling-for-on-chain-credit-risk}
\addcontentsline{toc}{subsection}{3.2 온체인 신용 리스크 예측 모델링}

그래프 이론적 방법과 병행하여, 연구자들은 특히 DeFi 대출 맥락에서 온체인 신용 점수 산출을 위한 데이터 기반 머신러닝 모델을 개발해 왔다. 이러한 모델은 전통적 신용 분석에서 영감을 받아 지갑 이력을 입력 특성으로 사용하고 부도 또는 연체 가능성을 예측한다. Ghosh et al.(2024)은 arXiv 프리프린트에서 공개 이벤트 로그로부터 지갑 활동을 특성화하고 지도 학습 모델을 훈련하여 차입자를 상환 및 청산 리스크 순으로 순위화하는 DeFi용 온체인 신용 리스크 점수 프레임워크를 제안한다. 이 틀에서 각 지갑의 활동은 신용 보고서와 유사하게 특성화된다: 특성에는 상환 이력, 부채 금액, 신용 이력 길이, 신규 신용, 신용 구성 등이 포함되며 모두 온체인 행동에서 파생된다. 머신러닝 분류기를 사용하여 주어진 온체인 대출 포지션이 가까운 미래에 연체(예: 담보 요건 미달 및 청산 직면)될 확률을 추정한다. 출력은 본질적으로 지갑 또는 포지션의 신용 리스크 점수로, 부도 확률을 나타낸다. Ghosh et al.은 이 접근이 차입자를 리스크 순으로 성공적으로 순위화할 수 있으며, 따라서 담보 비율만이 아닌 예측된 상환 확률에 의존하여 담보 부족 대출을 지원할 수 있음을 보고한다.


다른 머신러닝 접근도 마찬가지로 온체인 거래 및 계정 데이터를 신용도 예측 모델의 입력으로 취급한다. Hassija et al.(2020)은 전망 이론(prospect theory)을 활용하여 긍정적·부정적 신용 이벤트에 다르게 가중하는 블록체인 기반 신용 점수 모델을 제시했다. 그들의 모델은 스마트 계약으로 구현되어 새 대출 결과가 발생할 때마다 각 참여자의 신뢰 점수를 갱신한다.


전망 이론을 통합함으로써(예: 적시 상환보다 부도에 더 큰 페널티), 점수 메커니즘은 인간의 리스크 편향과 정렬되어 누가 부도할지 예측을 개선할 수 있다. 유사하게 Chakraborty et al.(2019)은 고객 신용 데이터를 온체인에 저장·평가하여 더 효율적인 금융 시스템 워크플로를 지원하는 블록체인 기반 신용 분석 프레임워크를 제안한다. 그들의 설계는 분산 원장 구조가 전통적 거래 이력과 함께 감사 가능한 신용 평가를 어떻게 지원할 수 있는지 보여준다.


\phantomsection
\subsection*{3.3 DeFi 신용 평가에 관한 산업 이니셔티브 및 학술적 관점}
\label{sec:industry-initiatives-and-non-scholarly-perspectives}
\addcontentsline{toc}{subsection}{3.3 DeFi 신용 평가에 관한 산업 이니셔티브 및 학술적 관점}

온체인 신용 점수 산출에 대한 산업 수요는 상업적 배포와 탈중앙화 대출 워크플로에 대한 학술적 분석을 모두 촉발했다. Moln\'{a}r et al.(2023)은 금융에서 블록체인 기반 비즈니스 프로세스 관리를 검토하고, 신용 및 청구 요청이 스마트 계약 워크플로를 통해 어떻게 재설계될 수 있는지 보여주며 DeFi 대출에서 투명하고 감사 가능한 리스크 신호의 필요성을 강조한다. Packin and Lev-Aretz(2024)는 하이브리드 온체인/오프체인 신용 아키텍처와 전통적 신용 데이터를 블록체인 네이티브 활동과 통합할 때의 거버넌스 과제를 분석하며, 평판 지표가 가명성 제약 하에서 해석 가능해야 함을 강조한다.


최근 학술 연구는 DeFi 신용 점수 산출을 순수 상업 기능이 아닌 측정 가능한 연구 문제로 다룬다. Shimpi(2024)는 전통적 신용 점수 구성요소를 스마트 계약 신용국 기능과 통합하는 이더리움 기반 대출 워크플로를 구현하여, 오프체인 신용 데이터가 온체인 대출 결정에 어떻게 연결될 수 있는지 보여준다. Vathsala V. and V. D. Sharma(2025)는 연합 학습과 블록체인 고정 데이터 출처를 결합한 DeFi 동적 신용 점수를 위한 보안 AI 아키텍처를 제안하며, 대부분의 배포 시스템이 명시적 승인 관계가 아닌 송금 이력 또는 대출 결과에 의존한다고 지적한다.


이러한 학술적 관점은 EndorseRank가 채우려는 현재 공백을 강조한다. 기존 구현은 과거 거래 행동(상환, 청산, 잔액) 또는 지도 학습 부도 예측에 크게 집중하지만, 지갑 간 ERC-20 토큰 승인 관계를 주요 그래프 엣지로 체계적으로 활용하지는 않는다. ERC-20 approve(또는 permit) 행위는 다른 주소가 소유자를 대신하여 토큰을 지출하도록 승인하며(Ethereum Improvement Proposals, 2015; Ethereum Improvement Proposals, 2020), 직접적인 가치 이전 밖에 있는 강한 신뢰 또는 금융 연결 신호이다. 현재 그래프 기반 점수 시스템은 이러한 승인 엣지를 대체로 무시하고 실현된 가치 이전에만 집중한다. EndorseRank는 ERC-20 approve/permit 엣지를 평판 그래프의 일급 링크로 통합하여 PageRank 계산에서 방향성 가중 엣지로 취급함으로써 이 연구 공백을 해결한다.


또한 EndorseRank는 시간적 효과 처리에서 독특하다. 많은 평판 시스템은 오래된 상호작용의 가중치를 낮추기 위해 시간 기반 감쇠를 적용하는 반면, EndorseRank는 approve/permit 엣지에 대한 시간 감쇠를 생략할 것을 제안한다. 근거는 승인이 취소될 때까지 유효하다는 것이다; 1년 전에 부여된 권한도 승인이 여전히 활성 상태이면 오늘날 동등하게 관련 있다. 오래 지속된 신뢰 링크에 지속적 가중치를 부여함으로써 EndorseRank는 짧은 기억 창이 잃어버릴 중요한 지속적 금융 관계 정보를 보존한다. 이는 수년 전 결제에 합리적으로 낮은 중요성을 부여할 수 있는 전형적 거래 기반 점수와 대조된다. 승인 맥락에서 시간 감쇠 부재는 지속적 리스크의 본질을 반영한다: B가 A의 자금 지출을 승인하는 한, A의 노출(그리고 따라서 평판적 연결)은 시간이 지나도 감소하지 않는다. 이 no-decay 접근은 독특하며, 대부분의 선행 연구는 이러한 엣지를 전혀 포함하지 않았거나 신뢰의 시간적 측면을 명시적으로 고려하지 않았다. ERC-20 승인 네트워크를 PageRank와 통합하고 시간에 따른 영향을 유지함으로써, EndorseRank는 가치 흐름과 Arbitrum One의 위임된 권한 망을 모두 포착하는 더 풍부하고 안정적인 온체인 평판 점수를 제공할 것으로 기대된다.


요약하면, 문헌은 단순 그래프 기반 평판 지표에서 정교한 머신러닝 온체인 신용 점수 모델로의 진행을 보여준다. PageRank 기반 기법은 신뢰와 영향력에 대한 투명하고 설명 가능한 네트워크 중심적 관점을 기여하고, ML 모델은 다양한 특성을 사용한 신용 리스크 순위 성능을 제공한다. 그러나 두 가지를 융합하는 공백이 남아 있다: ERC-20 승인 링크와 같은 의미 있는 그래프 특성을 설명 가능하면서 도메인에 정렬된 점수 시스템에 통합하는 것이다. EndorseRank는 이 공백을 메우려 한다. 네트워크 구조가 중요하다---누가 누구를 신뢰하는가(송금과 승인 모두를 통해)---는 통찰과, 이 구조가 PageRank식 알고리즘을 통해 정량화될 수 있다는 통찰에 기반한다. 본 학위논문의 실증 검증은 Arbitrum One의 GMX V2 무기한 스왑 포지션 결과를 사용한다: 청산 성공 횟수와 실현 이익 대리변수는 Do et al.(2023)에서 사용된 in-degree 및 in-value 대리변수와 유사한 관찰 가능한 거래 성공 측정치 역할을 한다. 오래된 신뢰 링크를 임의로 폐기하지 않음으로써 EndorseRank는 관계의 장기 맥락을 보존한다. 이 접근은 선행 연구의 한계를 직접 다루며 투명하고 조합 가능한 DeFi 리스크 평가에 대한 학술적 요구(Moln\'{a}r et al., 2023; Shimpi, 2024)와 정렬한다. Arbitrum One의 레버리지 DeFi 활동에 대한 리스크 평가를 개선할 수 있는 향상된 온체인 평판 점수를 약속하며, 탈중앙화 금융을 더 접근 가능하고 더 안전하게 만든다.
